\documentclass[11pt,a4paper]{article}

\usepackage[utf8]{inputenc}
\usepackage[T1]{fontenc}
\usepackage{amsmath}
\usepackage{makeidx}
\usepackage[colorlinks=true,linkcolor=blue,citecolor=blue,urlcolor=blue]{hyperref}

\makeindex

\title{A Reading List on Information, Computation and Concurrency}
\author{OverBranch Benchmark Corpus}
\date{}

\begin{document}

\maketitle

\begin{abstract}
  This short note collects a small reading list used as a workload for the
  bibliography and index paths of the OverBranch compiler benchmark. It
  exercises \textsc{BibTeX} and \textsc{MakeIndex} together with ordinary
  cross-references.
\end{abstract}

\section{Information}

The quantitative study of communication begins with Shannon's
definition of entropy~\cite{shannon1948}. The \index{entropy} entropy of a
discrete source is $H = -\sum_i p_i \log_2 p_i$, and it sets the limit on
lossless compression of that source.

\section{Computation}

Turing framed the question of machine intelligence in terms of an imitation
game~\cite{turing1950}. Modern algorithm textbooks such as
\cite{cormen2009} give the standard treatment of \index{complexity}
complexity classes and asymptotic analysis.

The structure of a multiprogrammed operating system was described by
Dijkstra~\cite{dijkstra1968}, whose work on \index{semaphore} semaphores
remains influential. Hoare later introduced communicating sequential
processes~\cite{hoare1978}, a model of \index{concurrency} concurrency based
on synchronous message passing.

\section{Concurrency and time}

Lamport analysed the ordering of events in distributed systems and introduced
logical clocks~\cite{lamport1978}. The \index{happens-before relation} happens-before
relation he defined is a partial order on events.

\section{Typesetting tools}

Knuth's description of the \TeX{} system~\cite{knuth1984} and Lamport's
guide to \LaTeX{}~\cite{lamport1994} are the canonical references for the
typesetting toolchain. Documents built with these tools use an
\index{index}index generated from \verb|\index| commands, as this one does.

\bibliographystyle{plain}
\bibliography{refs}

\appendix

\section{Index of terms}

The index below is generated from the \verb|\index| commands in the body.

\printindex

\end{document}
